\documentclass[10pt,a4paper]{amsart}
\usepackage[margin=19mm]{geometry}
\usepackage{amsmath,amssymb,mathtools}
\usepackage[scr=rsfs,cal=euler]{mathalfa}
\usepackage{xfrac}
\usepackage[unicode,hidelinks,bookmarks=false]{hyperref}
\newcommand{\ie}{i.e.\ }
\newcommand{\cf}{cf.\ }
\newcommand{\resp}{respectively\ }
\newcommand{\Subsection}{\subsection}
\providecommand{\nobreakdash}{\nobreak\hbox{-}}
\setlength{\emergencystretch}{2em}
\multlinegap0pt
\usepackage{bm}
\usepackage{bbm}
\usepackage{dsfont}
\usepackage{xspace}
\usepackage{cancel}
\usepackage{mathtools}
\usepackage[shortlabels]{enumitem}
\usepackage{amsthm}
\makeatletter
\@ifundefined{thm}{\newtheorem{thm}{Thm}}{}
\@ifundefined{cor}{\newtheorem{cor}[thm]{Corollary}}{}
\@ifundefined{definition}{\newtheorem{definition}[thm]{Definition}}{}
\@ifundefined{lemma}{\newtheorem{lemma}[thm]{Lemma}}{}
\@ifundefined{prop}{\newtheorem{prop}[thm]{Proposition}}{}
\makeatother
\newcommand\N{\mathbb{N}}
\newcommand\calL{\mathcal{L}}
\newcommand{\fr}{\mathfrak}
\title[Dimension theory for the asymptotic couple of the field of l]{Dimension theory for the asymptotic couple of the field of logarithmic transseries}
\author{Allen Gehret, Elliot Kaplan, Nigel Pynn-Coates}
\date{Source: arXiv:2507.01760v2 (CC BY 4.0)}
\begin{document}
\maketitle

\noindent\emph{Source and licence.} Dimension theory for the asymptotic couple of the field of logarithmic transseries. Version arXiv:2507.01760v2 (CC BY 4.0), distributed under the Creative Commons Attribution 4.0 International licence, as recorded in the arXiv licence metadata for this exact version and shown on the abstract page https://arxiv.org/abs/2507.01760v2. Full rights notice: licenses/arxiv-2507.01760-CC-BY-4.0.txt.\par\smallskip

\noindent\textit{Editorial scope.} Counted unit: the Prop whose source label is \texttt{d-minimality\_criterion\_ms} in the article (source0.tex lines 2052--2061, source label \texttt{d-minimality\_criterion\_ms}), delivered together with the article's own complete original proof of it (source0.tex lines 2165--2194, one \texttt{proof} environment of 30 lines). No mathematics is written, repaired or re-derived: statement and proof are the article's own text line for line. The proof is detached from the statement in the source: the article states the result at lines 2052--2061 and prints its proof at lines 2165--2194, and the proof environment's own header names the statement, so the association is the article's own. Required context retained uncounted: all\_terms\_preps\_ms (lines 2158--2160); discrete\_gluing\_lemma (lines 2444--2452); def:stronglydiscrete (lines 2457--2459). Every definition, display and label of those lines is retained unchanged. Omitted from this package and not redistributed: 1--2051, 2062--2157, 2161--2164, 2195--2443, 2453--2456, 2460--2859. Nothing inside a retained excerpt is omitted or altered: every retained body is delivered exactly as the article's lines are written, so no omission receipt of any kind accompanies this package. Citations: the retained excerpts carry no citation command at all, so no bibliography entry of the article is transcribed and no reference is redistributed. Reference adaptation: none was needed and none was made; every reference inside the delivered unit resolves inside it, so no unresolved reference or citation is produced. Printed slips kept literal: no printed word, symbol, hypothesis or number of the counted statement or of its proof was altered. Layout and class: the article's own class is \texttt{amsart} with packages including amsmath, amsthm, amssymb, enumerate, fullpage, xy, hyperref, url, color, tikz, enumitem, mathtools, bm, stackengine; the wrapper is the lane's pinned \texttt{amsart} editorial wrapper at 10pt on A4, which restates the theorem-like environments the retained text needs and adds the article's own macro definitions that the retained text needs (\texttt{\textbackslash N}, \texttt{\textbackslash calL}, \texttt{\textbackslash fr}), with extra wrapper packages (bm, bbm, dsfont, xspace, cancel, mathtools, enumitem). The delivered numbering is the unit's own. Assets: no retained span contains a figure environment, a graphics-inclusion command, a PDF-page inclusion, a table or a TikZ picture; no image file is copied into this package, the package declares no asset dependency and no image file is redistributed with it. Licence: version arXiv:2507.01760v2 of this article is distributed under the Creative Commons Attribution 4.0 International licence, as recorded in the arXiv licence metadata for this exact version and shown on the abstract page https://arxiv.org/abs/2507.01760v2. A full rights notice is delivered at licenses/arxiv-2507.01760-CC-BY-4.0.txt. Characters: every retained excerpt is pure ASCII, so the delivered engine, XeLaTeX with its default Latin Modern text font, reports no missing character.
\begin{cor}\label{all_terms_preps_ms}
If each $\fr{f}\in\fr{F}$  is locally constant off of a finite union of strongly discrete sets, then every $\calL(\fr{F})$-term has a (continuous) preparation.
\end{cor}

\begin{lemma}\label{discrete_gluing_lemma}
Suppose:
\begin{enumerate}
\item $(B_i)_{i\in I}$ is a disjoint family of open sets,
\item $(D_i)_{i\in I}$ is a family with $D_i \subseteq B_i$ for each $i \in I$,
\item there is $N \in \N$ such that each $D_i$ is a union of $N$ discrete sets.
\end{enumerate}
Then $\bigsqcup_{i\in I}D_i$ is a union of $N$ discrete sets. 
\end{lemma}

\begin{definition}\label{def:stronglydiscrete}
We say a set $D\subseteq X$ is \textbf{strongly discrete} if there exists a family $(V_d)_{d\in D}$ of pairwise disjoint open sets such that $V_d\cap D=\{d\}$ for each $d\in D$.
\end{definition}

\begin{prop}[d-minimality criterion]\label{d-minimality_criterion_ms}
Suppose the following conditions hold:
\begin{enumerate}[label=(\Alph*)]
\item\label{conditionA} For every $\calL(\fr{F})$-formula $\varphi(t)$ where $t$ is a unary  variable in the sort $s_0$, there exist an $\calL$-formula $\varphi^*(x_1,\ldots,x_n)$ and $\calL(\fr{F})$-terms $\tau_1(t),\ldots,\tau_n(t)$ such that:
\[
T(\fr{F}) \ \vdash \ \varphi(t)\leftrightarrow \varphi^*(\tau_1(t),\ldots,\tau_n(t)).
\]
\item\label{conditionB} Every new function symbol $\fr{f}\colon M_{s'}\to M_{s}$ in $\fr{F}$ is locally constant off of a finite union of strongly discrete sets; see Definition~\ref{def:stronglydiscrete}.\end{enumerate}
Then $T(\fr{F})$ is d-minimal at $s_0$.
\end{prop}

\begin{proof}[Proof of Proposition~\ref{d-minimality_criterion_ms}]
Let $D\subseteq M_0$ be $\calL(\fr{F})$-definable. By removing the interior of $D$ (which is open and $\calL(\fr{F})$-definable), we may assume that $D$ has empty interior. We will show that $D$ is a finite union of discrete sets. 
By \ref{conditionA}, there exist an $n$-ary $\calL$-definable relation $R(x_1,\ldots,x_n)$ and $\calL(\fr{F})$-terms $\tau_1(t),\ldots,\tau_n(t)$ such that $D$ is of the form:
\[
D \ = \ \{t\in M_0:R(\tau_1(t),\ldots,\tau_n(t))\}.
\]
Next, for $i=1,\ldots,n$, by \ref{conditionB} and Corollary~\ref{all_terms_preps_ms} we take preparations $(\bm{s}_{i},B_{i},f_{i})$ of each term $\tau_{i}$. Set $\bar{\bm{s}}\coloneqq (\bm{s}_1,\ldots,\bm{s}_n)$ and define  $B\subseteq M_{\bar{\bm{s}}}\times M_0$ by declaring for $x=(x_1,\ldots,x_n) \in M_{\bm{s}_1}\times \cdots \times M_{\bm{s}_n}$: 
\[
B_x \ \coloneqq  \ B_{1,x_1}\cap\cdots\cap B_{n,x_n}\ \subseteq M_0.
\]
Then each $B_x$ is open and 
\[
\textstyle M_0\setminus\bigcup_{x}B_x \ = \ \bigcup_{1\leq i\leq n} (M_0 \setminus \bigcup_{x_{i}\in M_{\bm{s}_i}}(B_{i})_{x_{i}})
\]
is a finite union of $\calL(\fr{F})$-definable discrete sets.
We have 
\begin{align*}
D \cap B_x\ &= \ \{ t \in B_x : R_{i}\big(f_{0}(x_{0},t),\dots,f_{n-1}(x_{n-1},t)\big)\ \} \\
&= \ B_x\cap 
\underbrace{\{ t \in M_0 : R_{i}\big(f_{0}(x_{0},t),\dots,f_{n-1}(x_{n-1},t)\big) \}}_{\eqqcolon C_x}.
\end{align*}
Suppose $z\in\operatorname{int}(C_x)$, so there is an open $U$ such that $z\in U\subseteq C_x$. Since $D\cap B_x$ does not have interior, it must be the case that $U\cap B_x=\varnothing$ (since $B_x$ is open), thus $z\not\in D\cap B_x$. In particular, $D\cap B_x=B_x\cap \operatorname{br}(C_x)$. By our assumption that $T$ is d-minimal at $s_0$, $\operatorname{br}(C_x)$ is a union of $N$ discrete sets, for some $N \in \N$ not depending on $x$, so $D \cap B_x$ is a union of $N$ discrete sets. 

By Lemma~\ref{discrete_gluing_lemma}, the set $D \cap \bigcup_{x \in M_{\bar{\bm{s}}}} B_x$ is a union of $N$ discrete sets.
Finally,
\[
\textstyle D\ =\ (D \cap \bigcup_{x \in M_{\bar{\bm{s}}}} B_x) \cup (D \cap (M_0 \setminus \bigcup_{x \in M_{\bar{\bm{s}}}} B_x))
\]
is a finite union of discrete sets.
\end{proof}

\end{document}
